\documentclass[a4paper,12pt]{report}

\usepackage[lmargin=2.00000cm,rmargin=2.0000cm,tmargin=5.5000cm,bmargin=2.500000cm,headheight=3.5cm]{geometry}        %Flexible and complete interface to document dimensions
\usepackage[utf8]{inputenc}
\usepackage[T1]{fontenc}
%\usepackage[latin1]{inputenc}
\usepackage{amsmath}
\usepackage{amsfonts}
\usepackage{amssymb}
\usepackage{lmodern}
\usepackage{float}
\usepackage{graphicx}

\usepackage[english,french]{babel}        %use for the below package `datetime'
\usepackage[babel=true]{csquotes}
\usepackage{datetime}        %Change format of `\today' with commands for current time
\renewcommand{\dateseparator}{-}
\newcommand{\headertoday}{\twodigit\day \dateseparator \twodigit\month  \dateseparator \the\year}

\usepackage{animate}
\usepackage{lastpage}
\usepackage{array}
\usepackage{lastpage}
\usepackage{multirow}
\usepackage{titling}
\usepackage{placeins}
\usepackage{media9}
\usepackage{eurosym}
\usepackage{pdflscape}
\usepackage{color}
\usepackage[table]{xcolor}
\definecolor{mon_bleu}{rgb}{0.137,0.466,0.741}
\definecolor{mon_vert}{rgb}{0.07843,0.4627,0.07843}
\definecolor{mon_rouge}{rgb}{0.62745,0.16078,0.27058}		
%\usepackage{subfigure}
\usepackage{subcaption}

\usepackage[ 
           hidelinks,
           colorlinks=true,
           linkcolor=blue,          % color of internal links (change box color with linkbordercolor)
    	   citecolor=blue,        % color of links to bibliography
    	   filecolor=blue,      % color of file links
    	   urlcolor=blue,           % color of external links
           pdfhighlight =/O]{hyperref}																		% dvipdfm package pour lien hypertextes dans pdf (pdfhighlight--> afficher la main dans le pdf, colorlinks--> colorlinks=true afficher les liens en couleurs)
  
\usepackage{blindtext}
\usepackage[final]{pdfpages}
\usepackage[french]{cleveref}																				% à placer après hyperref attention ne pas utiliser ":" dans les labels des équations avec French babel activé et cref

\usepackage{fancyhdr}
%%%%%%%%%%%%%%%%%%%%%%%%%%%%%%%%%%%%%%%%%%%%%%%%%%%%%%%%%%%%%%%%%%%%%%%%%%
%                         FORMAT PERSONNALISES                           %
% %%%%%%%%%%%%%%%%%%%%%%%%%%%%%%%%%%%%%%%%%%%%%%%%%%%%%%%%%%%%%%%%%%%%%%%%

\parindent=0pt        %leading space for paragraphs
\pagestyle{fancy}
\renewcommand{\arraystretch}{1.5}
\renewcommand{\headrulewidth}{0pt}
\fancyhead[CE,CO,LE,LO,RE,RO]{} %% clear out all headers
\fancyhead[C]{%
\begin{tabular}{|m{3.0cm}|m{10cm}|c@{}|}
\hline
\multirow{2}{*}{\includegraphics[scale=0.045]																			% LOGO LABO
{Logo-Symme.png}}& \centering \multirow{2}{*}{ \Large{\thetitle}}  &  \multirow{2}{*}{Page: \thepage ~/ \pageref{LastPage}}\\
& &  \\
\hline
%Etabli par: \newline \theauthor & \multicolumn{2}{l|}{Diffusion: interne}\\
%\hline
\multicolumn{2}{|l|}{Objet: \thetitleobject }& Date: \headertoday ~~\\ 
\hline
\end{tabular}
}
\renewcommand\footrulewidth{1pt}
\fancyfoot[L]{Document confidentiel} 																		% DOCUMENT CONFIDENTIEL
\fancyfoot[R]{Laboratoire SYMME}		 																	% NOM DU LABO

\newcolumntype{x}[1]{>{\centering\hspace{0pt}}p{#1}}
\setlength{\doublerulesep}{\arrayrulewidth} 


\usepackage{tikz}																							% pour les diagrammes
\usetikzlibrary{calc}
\usetikzlibrary{shapes.geometric,shapes.arrows,decorations.pathmorphing}
\usetikzlibrary{matrix,chains,scopes,positioning,shapes,arrows,fit}
\usetikzlibrary{calc,decorations.pathreplacing}
\usetikzlibrary{calc}
\usetikzlibrary{backgrounds,decorations.markings}

\setcounter{secnumdepth}{5}		%sous sous sous section



%%%%%%%%%%%%%%%%%%%%%%%%%%%%%%%%%%%%%%%%%%%%%%%%%%%%%%%%%%%%%%%%%%%%%%%%%%%%%%%%%%%%
%%-------------> PAGE DE GARDE INFO
%%%%%%%%%%%%%%%%%%%%%%%%%%%%%%%%%%%%%%%%%%%%%%%%%%%%%%%%%%%%%%%%%%%%%%%%%%%%%%%%%%%%

\author{Auteur}
\newcommand{\validator}{J. COLLOMB}
\title{Rapport d'avancement}
\selectlanguage{french}	
\date{\today}
\newcommand{\thetitleobject}{Comment gérer sa bibliographie ?}
\setcounter{tocdepth}{6}
\setcounter{secnumdepth}{6}


%%%%%%%%%%%%%%%%%%%%%%%%%%%%%%%%%%%%%%%%%%%%%%%%%%%%%%%%%%%%%%%%%%%%%%%%%%%%%%%%%%%%
%%-------------> DEBUT DU DOCUMENT 
%%%%%%%%%%%%%%%%%%%%%%%%%%%%%%%%%%%%%%%%%%%%%%%%%%%%%%%%%%%%%%%%%%%%%%%%%%%%%%%%%%%%

\begin{document}

\graphicspath{{Figures/}}

%%%%%%%%%%%%%%%%%%%%%%%%%%%%%%%%%%%%%%%%%%%%%%%%%%%%%%%%%%%%%%%%%%%%%%%%%%%%%%%%%%%%
%%-------------> PAGE DE GARDE 
%%%%%%%%%%%%%%%%%%%%%%%%%%%%%%%%%%%%%%%%%%%%%%%%%%%%%%%%%%%%%%%%%%%%%%%%%%%%%%%%%%%%
\begin{titlepage}
    \centering
    \vspace*{3cm}
    {\bfseries\Large
        \thetitle\\
        \thetitleobject\\
        --- \\
        Document confidentiel\\  																			% LIGNE "CONFIDENTIEL"
        \vskip2cm
    }
    \includegraphics[scale=0.65]{logos.png}  																% LOGOS DES PARTENAIRES
    \vfill
    \begin{center}
    \begin{table}[b]
    \begin{tabular}{x{.225\linewidth}|| x{.225\linewidth}|| x{.225\linewidth} || x{.225\linewidth} }
    \hline \hline
    \textbf{Date} & \textbf{Révision} & \textbf{Rédigé par} & \textbf{Validé par} \tabularnewline
    \hline \hline
    \thedate & Rev. A  &  \theauthor &  \validator \tabularnewline
    \hline
    ~ & ~ & ~ & ~ \tabularnewline
    \hline
    ~ & ~ & ~ & ~ \tabularnewline
    \hline
    ~ & ~ & ~ & ~\tabularnewline         
    \hline \hline                                                                   
    \end{tabular}
    \end{table}
    \end{center}
    \vfill
    \vfill
    \begin{tikzpicture}[remember picture,overlay]
    \node (label) at (8cm,20cm){
        \includegraphics[width=2.5cm]{Logo-Symme.png} 													% LOGO LABO
      };
    
	\draw[very thick]
		([yshift=-25pt,xshift=25pt]current page.north west)--
		([yshift=-25pt,xshift=-25pt]current page.north east)--
		([yshift=25pt,xshift=-25pt]current page.south east)--
		([yshift=25pt,xshift=25pt]current page.south west)--cycle;
    \end{tikzpicture}
\end{titlepage}


%\begin{titlepage}
%    \centering
%    
%    \vspace*{3cm}
%    {\bfseries\Large
%        \thetitle\\
%        \thetitleobject\\
%        \vskip2cm
%    }    
%    \vfill
%    \begin{center}
%    \begin{table}[b]
%    \begin{tabular}{x{.225\linewidth}|| x{.225\linewidth}|| x{.225\linewidth} || x{.225\linewidth} }
%    \hline \hline
%    \textbf{Date} & \textbf{R�vision} & \textbf{R�dig� par} & \textbf{Valid� par} \tabularnewline
%    \hline \hline
%    \thedate & Rev. A  &  \theauthor &  \validator \tabularnewline
%    \hline
%    ~ & ~ & ~ & ~ \tabularnewline
%    \hline
%    ~ & ~ & ~ & ~ \tabularnewline
%    \hline
%    ~ & ~ & ~ & ~\tabularnewline         
%    \hline \hline                                                                   
%    \end{tabular}
%    \end{table}
%    \end{center}
%    \vfill
%    \vfill
%    \begin{tikzpicture}[overlay,remember picture]
%    \node (label) at (6.7cm,19.5cm){
%        \includegraphics[width=4cm]{../Figures/logo_CT1.pdf} % also works with logo.pdf
%      };
%    
%    \draw [line width=1pt,rounded corners=7pt]
%        ($ (current page.north west) + (1.5cm,-1.5cm) $)
%        rectangle
%        ($ (current page.south east) + (-1.5cm,1.5cm) $);
%        
%%	\draw[very thick]
%%		([yshift=-25pt,xshift=25pt]current page.north west)--
%%		([yshift=-25pt,xshift=-25pt]current page.north east)--
%%		([yshift=25pt,xshift=-25pt]current page.south east)--
%%		([yshift=25pt,xshift=25pt]current page.south west)--cycle;
%    \end{tikzpicture}
%\end{titlepage}


%%%%%%%%%%%%%%%%%%%%%%%%%%%%%%%%%%%%%%%%%%%%%%%%%%%%%%%%%%%%%%%%%%%%%%%%%%%%%%%%%%%%
%%-------------> SOMMAIRE
%%%%%%%%%%%%%%%%%%%%%%%%%%%%%%%%%%%%%%%%%%%%%%%%%%%%%%%%%%%%%%%%%%%%%%%%%%%%%%%%%%%%

\renewcommand\contentsname{Sommaire}
\setcounter{chapter}{1}
\tableofcontents
%\listoffigures
%\listoftables


%%%%%%%%%%%%%%%%%%%%%%%%%%%%%%%%%%%%%%%%%%%%%%%%%%%%%%%%%%%%%%%%%%%%%%%%%%%%%%%%%%%%
%%-------------> CORPS DOCUMENT
%%%%%%%%%%%%%%%%%%%%%%%%%%%%%%%%%%%%%%%%%%%%%%%%%%%%%%%%%%%%%%%%%%%%%%%%%%%%%%%%%%%%

%-----------------------------------------------------------------------------------
%-----------------------------------------------------------------------------------
\newpage

\section{Bibliographie ?}
\subsection{Pour quelle raison ?}
\textit{Le travail de recherche et l'écriture d'un texte scientifique (rapport, article, thèse,…) exigent une recherche d'informations approfondie qui prend directement appui sur les travaux antérieurs. L'information choisie et exploitée permet de développer une réflexion personnelle et chaque document, retenu et analysé, contribue à la crédibilité scientifique du travail présenté. \\}

\textit{Afin de faciliter la réflexion et le travail de recherche des lecteurs, qui à leur tour vont vouloir croiser leurs informations, il est indispensable de référencer correctement les travaux cités dans le texte en rédigeant une partie intitulée « Bibliographie » ou « Références bibliographiques ». \\}

\textit{La bibliographie d'un document permet de connaître : (i)les travaux qui ont été utilisés pour le travail de recherche et la rédaction ; (ii) l’état de la littérature sur un sujet pendant une période déterminée ;(iii) les auteurs, titres de revue, sites web… spécialisés dans un domaine.} \\

Source : \href{http://www.ajar-online.fr/thesememoire-2-recherche-bibliographique/}{AJAR Paris}    \\

Quelques endroits pour effectuer sa recherche bibliographique :
\begin{itemize}
\item \href{https://scholar.google.fr/}{Google Scholar};
\item \href{https://www-sciencedirect-com.camphrier-1.grenet.fr/}{Science Direct};
\item \href{https://link-springer-com.camphrier-1.grenet.fr/}{Springer};
\item \href{https://www-techniques-ingenieur-fr.camphrier-1.grenet.fr/}{Techniques de l'ingénieur}
\item \href{https://hal.archives-ouvertes.fr/}{HAL}.
\end{itemize}

\ \\

\begin{figure}[hbtp]
	\centering
	\def\svgwidth{1\columnwidth}
	\fontsize{10pt}{10pt}\selectfont\input{Figures/image_bibliographie.pdf_tex}
	\caption{Acquisition de connaissances}
	\label{figure_bouquins}
\end{figure}



\subsection{Comment faire ?}
L'utilisation d'un logiciel de gestion bibliographique est plus que recommandé dans le cadre de la thèse. Ces outils permettent généralement d'organiser sa bibliographie (en dossier par exemple), d'obtenir l'ensemble des informations nécessaires pour les citations dans les rapports scientifiques, d'exporter et d'intégrer aisément la bibliographie à un rapport\ldots \\

Il existe différents logiciels, comme par exemple : 
\begin{itemize}
\item \href{https://www.mendeley.com/}{Mendeley}; 
\item \href{https://www.zotero.org/}{Zotero};
\item \ldots 
\end{itemize}

Un exemple de l'interface de Mendeley est visible Figure~\ref{figure_mendeley}.

\begin{figure}[hbtp]
	\centering
	\def\svgwidth{1\columnwidth}
	\fontsize{10pt}{10pt}\selectfont\input{Figures/mendeley.pdf_tex}
	\caption{Interface de Mendeley}
	\label{figure_mendeley}
\end{figure}


\subsection{Faire une citation sous \LaTeX}
4 étapes sont nécessaires :
\begin{enumerate}
\item Exporter sa bibliographie au format \verb|.bib|;
\item Effectuer les citations;
\item Ajouter la bibliographie dans \LaTeX; 
\item Compiler.
\end{enumerate}

Les citations sont réalisées à l'aide de la commande : \verb|\cite{clé_de_citation}|\footnote{clé de citation : information renseignée dans le logiciel de gestion de bibliographie}. \\ 

L'ajout de la bibliographie s'effectue par les commandes : \verb|\bibliographystyle{ieeetr}| et \verb|\bibliography{nom_de_la_bibliographie_exportée}|. \\

La compilation consiste à compiler en pdf puis bibtex puis pdf (x2). Cela est fait automatiquement par la compilation rapide de TexMaker.



\section{Exemple}

Voici un super exemple dans lequel je peux faire une première citation \cite{collomb2017}, puis une deuxième \cite{collomb2018} et enfin une multi-citation \cite{collomb2018a, collomb2017a}.


%-----------------------------------------------------------------------------------
%-----------------------------------------------------------------------------------

\bibliographystyle{ieeetr}
\bibliography{bibliographie}


%%%%%%%%%%%%%%%%%%%%%%%%%%%%%%%%%%%%%%%%%%%%%%%%%%%%%%%%%%%%%%%%%%%%%%%%%%%%%%%%%%%%
%%-------------> FIN DU DOCUMENT
%%%%%%%%%%%%%%%%%%%%%%%%%%%%%%%%%%%%%%%%%%%%%%%%%%%%%%%%%%%%%%%%%%%%%%%%%%%%%%%%%%%%

\end{document}